% =========================================================
%  sections/07-pengalaman-belajar.tex
%  Bagian 7 — Pengalaman Belajar Mahasiswa
% =========================================================

\section{Pengalaman Belajar Mahasiswa}

Pengalaman belajar mahasiswa diwujudkan dalam rangkaian aktivitas mingguan yang
secara langsung mengarah pada pencapaian Sub-CPMK. Setiap pertemuan memiliki
keluaran yang dapat diamati: berkas kode, level yang dapat dimainkan, atau
laporan evaluasi. Tabel~\ref{tab:mingguan} menyajikan rencana pembelajaran
mingguan untuk 16 minggu, sedangkan spesifikasi tugas dan proyek diuraikan
pada bagian setelahnya.

\subsection{Rencana Pembelajaran Mingguan}

\begin{landscape}
{\setlength{\tabcolsep}{3pt}\footnotesize%
\begin{longtable}{@{}L{1.0cm} L{2.0cm} L{4.2cm} L{3.0cm} L{7.8cm} L{3.4cm} L{2.0cm}@{}}
\caption{Rencana pembelajaran mingguan mata kuliah \NamaMK}
\label{tab:mingguan}\\
\toprule
\headrow \textbf{Mg} & \textbf{Sub-CPMK} & \textbf{Pokok Bahasan} & \textbf{Metode \& Bentuk} & \textbf{Pengalaman Belajar Mahasiswa} & \textbf{Asesmen} & \textbf{Ref} \\
\midrule
\endfirsthead
\toprule
\headrow \textbf{Mg} & \textbf{Sub-CPMK} & \textbf{Pokok Bahasan} & \textbf{Metode \& Bentuk} & \textbf{Pengalaman Belajar Mahasiswa} & \textbf{Asesmen} & \textbf{Ref} \\
\midrule
\endhead
\midrule
\multicolumn{7}{r}{\footnotesize\itshape bersambung ke halaman berikutnya} \\
\endfoot
\bottomrule
\endlastfoot

1 & Orientasi (menuju 1.1) &
Ruang lingkup mata kuliah dan kontrak belajar; keterkaitan \CPL{03}--\CPL{04}, CPMK, dan Sub-CPMK; \istilah{game} sebagai simulasi interaktif \istilah{real-time}; peran algoritma dan bahasa pemrograman; instalasi Python 3, Pygame, dan editor; struktur folder proyek. &
Kuliah interaktif, demonstrasi, praktik terbimbing &
Menyiapkan lingkungan kerja (Python 3, Pygame, editor, Git); menjalankan skrip pertama yang membuka jendela Pygame dan menampilkan versi Pygame; menyusun struktur folder proyek dan mengunggahnya ke repositori; mendiskusikan kontrak kuliah dan capaian pembelajaran. &
Tugas Praktikum 01 (unjuk kerja) & [R1], [R2], [R20], [R21] \\
\addlinespace

2 & 1.1 &
Siklus input $\rightarrow$ \istilah{update} $\rightarrow$ \istilah{render}; \istilah{event handling} (\code{pygame.event.get}, \code{QUIT}); \code{pygame.time.Clock}; pembatasan FPS; konsep \istilah{fixed} dan \istilah{variable timestep}. &
Kuliah, \istilah{live coding}, praktik terbimbing &
Menulis \istilah{game loop} lengkap secara mandiri; menangkap \istilah{event} keluar dan tombol; membatasi laju bingkai dengan \code{Clock}; menampilkan nilai FPS pada judul jendela; mengubah batas FPS dan mencatat pengaruhnya terhadap beban prosesor. &
Tugas Praktikum 02; observasi kinerja \istilah{loop} & [R3], [R5], [R6], [R22] \\
\addlinespace

3 & 1.2 &
\code{pygame.Vector2}; posisi dan kecepatan; \code{pygame.key.get\_pressed}; \istilah{delta time} (\code{dt = clock.tick(FPS) / 1000}); sinkronisasi posisi \istilah{float} dengan \code{Rect} bilangan bulat. &
Kuliah, \istilah{live coding}, praktik terbimbing &
Mengimplementasikan pergerakan karakter dengan WASD atau tombol panah menggunakan \code{Vector2} dan \istilah{delta time}; membandingkan kecepatan objek pada 30 FPS dan 60 FPS; memperbaiki ketidaksesuaian antara posisi bertipe pecahan dan \code{Rect} bertipe bilangan bulat. &
Tugas Praktikum 03 & [R3], [R32], [R1] \\
\addlinespace

4 & 1.3 &
Percepatan gravitasi; \istilah{jump impulse}; pantulan dinding dan lantai; integrasi Euler sederhana; batas lantai. &
Kuliah, praktik terbimbing, latihan terbuka &
Menerapkan gravitasi, impuls lompatan, dan pantulan menggunakan integrasi Euler; memastikan lompatan hanya aktif ketika objek menyentuh tanah; menguji perilaku objek pada nilai gravitasi dan koefisien pantulan yang berbeda. &
Tugas Praktikum 04 & [R7], [R22], [R1] \\
\addlinespace

5 & 1.4 &
\code{Sprite} dan \code{Group}; \code{spritecollide} dan \code{groupcollide}; AABB dengan \code{Rect}; pengantar \istilah{pixel-perfect} dengan \code{mask}; resolusi \istilah{collision} sumbu demi sumbu. &
Kuliah, praktik terbimbing, \istilah{problem-based learning} &
Mengorganisasi objek permainan ke dalam \code{Sprite} dan \code{Group}; memilih serta menerapkan metode \istilah{collision} yang sesuai; menyelesaikan tabrakan per sumbu sehingga objek tidak menembus dinding; membangun mekanik pengumpulan item dan penghindaran rintangan dengan sistem skor. &
Tugas Praktikum 05; \textbf{Kuis 1} (Materi 1--4) & [R4], [R23], [R21] \\
\addlinespace

6 & 2.1 &
Persepsi jarak dan arah; normalisasi vektor menuju pemain; batasan kecepatan; perbedaan AI reaktif dan AI berbasis perencanaan. &
Kuliah, \istilah{live coding}, praktik terbimbing &
Menghitung vektor arah NPC menuju pemain dan menormalisasinya; membatasi besar kecepatan agar gerak konsisten; memvisualisasikan arah pengejaran dengan garis; membandingkan perilaku AI reaktif dengan AI yang mempertimbangkan perencanaan. &
Tugas Praktikum 06 & [R1], [R2], [R19], [R33] \\
\addlinespace

7 & 2.2 &
\istilah{State}, \istilah{event}, \istilah{transition}; diagram idle--patrol--chase--attack/flee; implementasi sebagai kelas atau fungsi Python; kondisi jarak dan \istilah{cooldown}. &
Kuliah, praktik terbimbing, \istilah{problem-based learning} &
Merancang diagram \istilah{state}--transisi NPC; mengimplementasikan FSM sebagai kelas atau fungsi Python; menambahkan transisi berbasis jarak dan \istilah{cooldown}; menguji setiap \istilah{state} dan transisinya secara terpisah. &
Tugas Praktikum 07 & [R9], [R25], [R27] \\
\addlinespace

8 & 2.3 &
Representasi grid dan graf; biaya $g$, heuristik $h$, dan $f = g + h$; jarak Manhattan dan Euclidean; antrean prioritas \code{heapq}; pemulihan jalur. &
Kuliah, \istilah{live coding}, praktik terbimbing &
Membangun representasi grid dan graf berbiaya dari peta berhalangan; mengimplementasikan A* dengan \code{heapq} dan heuristik Manhattan atau Euclidean; memulihkan jalur dan memvisualisasikannya pada grid Pygame; menganalisis pengaruh heuristik terhadap jumlah simpul yang diperiksa. &
Tugas Praktikum 08; \textbf{Kuis 2} (Materi 6--7) & [R8], [R10], [R11], [R34] \\
\addlinespace

9 & 2.4 &
FSM memanggil A* pada \istilah{state chase}; \istilah{repath} berkala; batasan performa; skenario uji (labirin dan rintangan dinamis sederhana). &
\istilah{Project-based learning}, kerja kelompok, presentasi &
Mengintegrasikan FSM dengan A*: \istilah{state chase} memanggil A* dan melakukan \istilah{repath} secara berkala; menyusun skenario uji berupa labirin dan rintangan dinamis sederhana; melaksanakan \istilah{playtest} antar kelompok; menyusun laporan singkat hasil \istilah{playtest} beserta perbaikan yang dilakukan. &
Presentasi dan laporan \textbf{Mini Proyek AI} & [R8], [R9], [R26], [R25] \\
\addlinespace

10 & 1.1--2.4 &
Asesmen tengah semester atas \CPMK{1} dan \CPMK{2}. &
Ujian tertulis dan uji kode &
Mengerjakan soal konsep \istilah{game loop}, manajemen waktu, vektor, fisika sederhana, \istilah{collision}, FSM, dan A*; menyelesaikan tugas perbaikan \istilah{bug} (\istilah{debugging}) pada kode Pygame yang disediakan dalam batas waktu tertentu. &
\textbf{UTS} & --- \\
\addlinespace

11 & 3.1 &
Unity Hub dan Editor; \istilah{Hierarchy}, \istilah{Inspector}, \istilah{Project window}; \istilah{scene}; \istilah{GameObject} dan komponen; perbandingan konsep Unity dengan Pygame. &
Kuliah, demonstrasi, praktik terbimbing &
Menavigasi antarmuka Unity; membuat \istilah{project} 2D; menyusun \istilah{scene} berisi kamera dan objek primitif; menyimpan \istilah{scene} dan mengelola aset; membandingkan istilah Unity dengan konsep Pygame yang telah dipelajari. &
Tugas Praktikum 09 & [R12], [R13], [R15], [R16] \\
\addlinespace

12 & 3.2 &
\istilah{Tilemap} dan \istilah{greyboxing}; \istilah{Collider2D/3D}; \istilah{Rigidbody}; \istilah{script} gerak pada \code{Update} dan \code{FixedUpdate}; \istilah{trigger volume}; desain jalur pemain. &
Demonstrasi, praktik terbimbing, \istilah{problem-based learning} &
Menyusun level dengan \istilah{tilemap} dan \istilah{greyboxing}; mengonfigurasi \istilah{Collider} dan \istilah{Rigidbody}; menulis \istilah{script} kontrol pemain; menambahkan \istilah{trigger volume}; menganalisis dan memperbaiki masalah fisika seperti \istilah{jitter}, objek menembus \istilah{collider}, atau jatuh tanpa henti. &
Tugas Praktikum 10; penilaian level & [R12], [R13], [R18], [R35] \\
\addlinespace

13 & 4.1 &
\istilah{Pipeline} impor Unity; \istilah{sprite sheet} dan \istilah{Animator}; \istilah{Canvas} UI (skor dan nyawa); penggantian \istilah{placeholder greybox} dengan aset sesungguhnya. &
Demonstrasi, praktik terbimbing, kerja kelompok &
Mengimpor \istilah{sprite} dan mengatur \istilah{importer setting}; menyiapkan \istilah{Animator} beserta transisi animasi idle dan run; membangun UI \istilah{Canvas} untuk skor dan nyawa; mengganti \istilah{placeholder} dengan aset visual; membandingkan \istilah{scene} sebelum dan sesudah integrasi aset. &
Tugas Praktikum 11 & [R12], [R14], [R28], [R29] \\
\addlinespace

14 & 4.2 &
\code{AudioClip} dan \code{AudioSource}; BGM dan SFX; \istilah{trigger} pada \istilah{pickup}, tabrakan, serta kondisi menang dan kalah; partikel sederhana (opsional). &
Demonstrasi, praktik terbimbing, kerja kelompok &
Menambahkan \code{AudioSource} dan \code{AudioClip} untuk BGM serta SFX; memicu efek suara pada peristiwa lompat, mengambil item, dan \istilah{game over}; menambahkan partikel sederhana (opsional); menyesuaikan \istilah{timing} dan level volume agar umpan balik terasa tepat. &
Tugas Praktikum 12 & [R14], [R24], [R30], [R31] \\
\addlinespace

15 & 4.3 &
\istilah{Checklist} fungsional; \istilah{playtest} antar kelompok; \istilah{log bug}; kesesuaian GDD dengan \istilah{build}; presentasi demo; refleksi pilihan Python dan Unity. &
\istilah{Project-based learning}, \istilah{playtest}, \istilah{peer review}, presentasi &
Menyusun \istilah{checklist} fungsional dan lembar \istilah{playtest}; melaksanakan \istilah{playtest} antar kelompok; mencatat \istilah{bug} dan masalah integrasi; memfinalisasi \istilah{build}; menyajikan demo 5--10 menit; menyusun laporan perbaikan dan refleksi pemilihan teknologi (Python/Pygame atau Unity). &
Presentasi \textbf{Proyek Akhir}; laporan perbaikan & [R12], [R16], [R18], [R30] \\
\addlinespace

16 & 4.3 &
Asesmen akhir semester atas \CPMK{3} dan \CPMK{4}. &
Uji fungsional, presentasi, dan penilaian dokumen &
Menyerahkan \istilah{build} dan \istilah{project source} prototipe \istilah{game}; menjalani uji fungsional oleh dosen menggunakan \istilah{checklist} penilaian; menyerahkan laporan pengembangan serta lembar refleksi pembelajaran. &
\textbf{UAS} & --- \\

\end{longtable}}
\tabsource{Keterangan: Mg = minggu ke-; Ref = nomor rujukan pada Bagian~\ref{sec:referensi}; UTS = Ujian Tengah Semester; UAS = Ujian Akhir Semester; GDD = \istilah{Game Design Document}.}
\end{landscape}

\subsection{Tugas dan Proyek}

\subsubsection{Tugas Praktikum (Portofolio Individu)}

Tugas praktikum dikerjakan setiap minggu secara individual dan dikumpulkan
melalui repositori Git pribadi mahasiswa. Setiap tugas memuat kode program yang
dapat dijalankan, tangkapan layar hasil, dan catatan singkat kendala yang
dihadapi beserta cara penyelesaiannya. Dosen menilai enam sampel portofolio
secara acak sebagai dasar nilai portofolio, sehingga mahasiswa terdorong
mengerjakan seluruh rangkaian tugas secara berkelanjutan.

\subsubsection{Kuis}

Kuis dilaksanakan dua kali dalam bentuk tes singkat (15--20 menit) untuk
memantau pemahaman konsep secara formatif.

\begin{itemize}
  \item \textbf{Kuis 1} (Minggu 5) mencakup Materi 1--4: struktur \istilah{game loop}, manajemen waktu, vektor, dan fisika sederhana.
  \item \textbf{Kuis 2} (Minggu 8) mencakup Materi 6--7: perilaku \istilah{chase} dan \istilah{Finite State Machine}.
\end{itemize}

\subsubsection{Mini Proyek AI (Kelompok, Materi 6--9)}

Mini proyek AI dikerjakan dalam kelompok beranggotakan 2--3 mahasiswa dan
dikumpulkan pada Minggu 9.

\begin{table}[htbp]
\centering
\caption{Spesifikasi mini proyek AI}
\label{tab:miniproyek}
\begin{tabularx}{\textwidth}{@{}L{3.4cm} Y@{}}
\toprule
\headrow \textbf{Aspek} & \textbf{Ketentuan} \\
\midrule
Bentuk keluaran &
Kode program Python/Pygame, ringkasan desain (\istilah{Game Design Document} singkat), dan laporan \istilah{playtest} 2--3 halaman. \\
\addlinespace
Isi wajib &
Satu atau lebih NPC yang menerapkan FSM (minimal \istilah{idle}, \istilah{patrol}, dan \istilah{chase}) dan menggunakan A* pada peta berhalangan saat berada pada \istilah{state chase}. \\
\addlinespace
Pengujian &
Menyertakan minimal tiga skenario uji (ruang terbuka, labirin, dan rintangan dinamis sederhana) beserta hasil pengamatan. \\
\addlinespace
Presentasi &
Demonstrasi 10 menit pada Minggu 9, dilanjutkan tanya jawab dan umpan balik dari kelompok lain. \\
\bottomrule
\end{tabularx}
\end{table}

\subsubsection{Proyek Akhir (Kelompok, Materi 10--15)}

Proyek akhir berupa prototipe \istilah{game} yang dibangun dengan Unity dan
dikumpulkan pada Minggu 16.

\begin{table}[htbp]
\centering
\caption{Spesifikasi proyek akhir}
\label{tab:proyekakhir}
\begin{tabularx}{\textwidth}{@{}L{3.4cm} Y@{}}
\toprule
\headrow \textbf{Aspek} & \textbf{Ketentuan} \\
\midrule
Bentuk keluaran &
\istilah{Build} yang dapat dijalankan, \istilah{project source}, ringkasan desain, laporan pengembangan, dan lembar \istilah{playtest} yang telah diisi penguji. \\
\addlinespace
Isi wajib &
Minimal satu level yang dapat dimainkan, kontrol pemain, mekanik inti, aset visual dengan animasi, elemen UI, serta aset audio yang terhubung pada peristiwa \istilah{gameplay}. \\
\addlinespace
Pengujian &
\istilah{Playtest} antar kelompok pada Minggu 15 menggunakan lembar penilaian terstruktur; seluruh \istilah{bug} yang ditemukan dicatat pada \istilah{log bug} dan diberi status penanganan. \\
\addlinespace
Presentasi &
Demo 5--10 menit pada Minggu 15 dan uji fungsional oleh dosen pada Minggu 16. \\
\bottomrule
\end{tabularx}
\end{table}

\subsubsection{Lembar Refleksi Pembelajaran}

Pada akhir semester setiap mahasiswa mengisi lembar refleksi yang memuat:
konsep yang paling sulit dipahami dan cara mengatasinya, perbandingan
pengalaman bekerja dengan Python/Pygame dan Unity, serta rencana pengembangan
kompetensi lanjutan. Lembar ini menjadi masukan perbaikan perkuliahan pada
semester berikutnya.
